% Notes of the Montreal seminar, 1970 (cote 129, typed part): archive cote 129, site page 58 (PDF page 59).
% typescript. Draft compared: none.
% AI-transcribed by Claude Opus 5.5 (claude-opus-5-5, adaptive thinking, effort high), Message Batches API, batch mtl129_full (msgbatch_01AGnDKHwy6ATCM7ehUa6263), 2026-11-13; not yet checked by eye.
\scanpage{129/58}
\mspage{III-3}

canoniques de $\Lambda_n$-modules constants qui donnent, par produit tensoriel
avec $G(n)$ les diagrammes suivants:

\diagram{
\begin{tikzcd}
p^i\Lambda_n \arrow[r, hook, "\varepsilon"] & \Lambda_n \\
\Lambda_n/p^{n-i}\Lambda_n \arrow[u, "p^i", "\wr"'] \arrow[ur, "(p^i)"'] &
\end{tikzcd}
\qquad
\begin{tikzcd}
p^i\Lambda_n\otimes G(n) \arrow[rr, "\varepsilon\otimes 1"] & & G(n) \\
\Lambda_n/p^i\Lambda_n\otimes G(n) = G(n)/p^iG(n) \arrow[u, equal, "\wr"'] \arrow[urr, "(p^i)"'] & &
\end{tikzcd}
}

les seuls idéaux de $\Lambda_n$ étant les $p^i\Lambda_n$, $G(n)$ est plat si et seulement
si $\varepsilon\otimes 1$ est injectif pour tout $i$, ou encore si et seulement si $(p^i)$
est injectif, ce qui équivaut à (c).

(c) $\Rightarrow$ (a) Les relations analogues à (c) sont alors vérifiées par
tous les $G(i)$ qui sont donc plats sur $\Lambda_i$, d'après ce qui précède.

(d) $\Rightarrow$ (e) Montrons que si $G(i) = p^{n-i}G(n)$, pour un certain $i$,
$1 \leq i \leq n-1$ alors on a $G(i+1) = p^{n-i-1}G(n)$. Cela se voit directement sur
le diagramme suivant où les trois lignes sont des suites exactes.

\diagram{
\begin{tikzcd}
0 \arrow[r] & G(n-1) \arrow[r, hook] \arrow[d, "p^{n-i-1}"] & G(n) \arrow[r, "p^{n-1}"] \arrow[d, "p^{n-i-1}"] & p^{n-1}G = p^{i-1}G(i) \arrow[r] \arrow[d, equal] & 0 \\
0 \arrow[r] & G(i) \arrow[r] \arrow[d] & G(i+1) \arrow[r, "p^i"] \arrow[d] & p^iG(i+1) = p^{i-1}G(i) \arrow[r] \arrow[d] & 0 \\
 & 0 \arrow[r, twoheadrightarrow] \arrow[d] & G(i+1)/p^{n-i-1}G(n) \arrow[r] \arrow[d] & 0 \arrow[r] \arrow[d] & 0 \\
 & 0 & 0 & 0 &
\end{tikzcd}
}

(e) $\Leftrightarrow$ (f) Soit $\lambda^i : p^iG(n)/p^{i+1}G(n) \rightarrow p^{i+1}G(n)/p^{i+2}G(n)$, $0 \leq i \leq n-2$

\mspage{58}
